\PassOptionsToPackage{numbers,sort&compress}{natbib}
\documentclass{article}

% paper_public.tex defines \PublicVersion before loading this file.
\providecommand{\PublicVersion}{0}
\ifnum\PublicVersion=1\relax
  \usepackage[preprint]{neurips_2026_vericode}
\else
  \usepackage{neurips_2026_vericode}
\fi
\workshoptitle{AI for Verifiable Coding}

\usepackage{amsmath,amssymb}
\usepackage{microtype}
\usepackage{booktabs,tabularx,array,longtable}
\usepackage{float}
\usepackage{graphicx}
\usepackage{xcolor}
\usepackage{listings}
\usepackage{listingsutf8}
\usepackage{tikz}
\usepackage[colorlinks=true,allcolors=blue]{hyperref}
\usepackage[nameinlink]{cleveref}
\usepackage[T1]{fontenc}
\usepackage[utf8]{inputenc}

% The review style loads lineno only for the anonymous submission build.
% Keep the historical theorem display usable in the public preprint as well.
\ifnum\PublicVersion=1\relax
  \newenvironment{nolinenumbers}{}{}
\fi

\newcommand{\Lean}{\textsc{Lean}}
\newcommand{\DK}{Davis--Kahan}
\input{generated/practitioner_account_macros.tex}
\input{generated/evidence_macros.tex}

% Set \PublicVersion to 1 in paper_public.tex for author names and public URLs.
\newif\ifshowartifacturl
\ifnum\PublicVersion=1\relax
  \showartifacturltrue
\else
  \showartifacturlfalse
\fi
\newcommand{\artifacturl}{https://github.com/AIQ-Kitware/aiq-dkps-formalization}

\definecolor{leankeyword}{HTML}{243B6B}
\definecolor{leantactic}{HTML}{70477A}
\definecolor{leantype}{HTML}{1F6473}
\definecolor{leancomment}{HTML}{39734D}
\definecolor{leanstring}{HTML}{8A4B2A}
\definecolor{leannumber}{HTML}{777777}
\definecolor{leanbackground}{HTML}{F7F7F3}
\lstdefinelanguage{lean}{%
  sensitive=true,
  morecomment=[l]{--},
  morecomment=[s]{/-}{-/},
  morestring=[b]",
  morekeywords=[1]{import,open,namespace,section,end,noncomputable,variable,variables,
    universe,universes,set_option,attribute,axiom,opaque,theorem,lemma,def,abbrev,
    example,instance,structure,class,inductive,notation,local,deriving,extends,
    by,where,let,have,show,fun,match,with,if,then,else,do,return,in,forall},
  morekeywords=[2]{simp,simpa,rw,erw,dsimp,unfold,ring,ring_nf,nlinarith,linarith,
    positivity,norm_num,field_simp,omega,exact,refine,apply,assumption,intro,intros,
    ext,funext,cases,rcases,obtain,constructor,left,right,classical,by_contra,
    contrapose,convert,using,suffices,calc},
  morekeywords=[3]{Prop,Type,Nat,Int,Real,Fin,Set,Submodule,Matrix,EuclideanSpace,
    LinearMap,ContinuousLinearMap,SymmetricNormingFunction,IsSelfAdjoint,Reduces},
}
\lstdefinestyle{leanpaper}{%
  language=lean,
  columns=fullflexible,
  backgroundcolor=\color{leanbackground},
  basicstyle=\ttfamily\footnotesize,
  commentstyle=\color{leancomment}\itshape,
  keywordstyle=[1]\color{leankeyword}\bfseries,
  keywordstyle=[2]\color{leantactic},
  keywordstyle=[3]\color{leantype},
  stringstyle=\color{leanstring},
  numberstyle=\tiny\color{leannumber},
  breakatwhitespace=false,
  breaklines=true,
  captionpos=b,
  keepspaces=true,
  numbers=none,
  showspaces=false,
  showstringspaces=false,
  showtabs=false,
  tabsize=2,
  aboveskip=0.6em,
  belowskip=0.6em,
  % This deliberately follows the MCC paper's \Real/\root pattern: the
  % listings input contains ASCII stand-ins and literate= does the typesetting.
  % IMPORTANT: the left-hand sides below are ASCII presentation sentinels,
  % not literal Lean Unicode. scripts/build_evidence.py writes the exact
  % extracted Lean signature to generated/historical_sin_theta_gap_mismatch_exact.lean,
  % then writes a separate listings input in which only display-sensitive Lean
  % notation is replaced by these sentinels. listings resolves each sentinel to
  % LaTeX here. Keep this indirection: it avoids Unicode decoding/editability
  % problems in listings/Overleaf while preserving an exact Lean sidecar for
  % audit. Do not replace the sentinel keys below with literal Lean glyphs.
  literate=
    {\\LeanLitCLMapC}{{$\mathrel{\to}\!\mathtt{L}[\mathbb{C}]$}}5
    {\\LeanLitCLMapK}{{$\mathrel{\to}\!\mathtt{L}[\mathbb{K}]$}}5
    {\\LeanLitComplex}{{$\mathbb{C}$}}1
    {\\LeanLitScalar}{{$\mathbb{K}$}}1
    {\\LeanLitReal}{{$\mathbb{R}$}}1
    {\\LeanLitLe}{{$\le$}}1
    {\\LeanLitSubsetEq}{{$\subseteq$}}1
    {\\LeanLitForall}{{$\forall$}}1
    {\\LeanLitMem}{{$\in$}}1
    {\\LeanLitLAngle}{{$\langle$}}1
    {\\LeanLitRAngle}{{$\rangle$}}1
    {\\LeanLitOrth}{{$^{\perp}$}}1
    {\\LeanLitDisj}{{$\lor$}}1
    {\\LeanLitConj}{{$\land$}}1
    {\\LeanLitPMapC}{{$\mathrel{\to}\!\mathtt{L}_{\rm p}[\mathbb{C}]$}}6
    {\\LeanLitPMapK}{{$\mathrel{\to}\!\mathtt{L}_{\rm p}[\mathbb{K}]$}}6
    {\\LeanLitArrow}{{$\to$}}1
    {\\LeanLitLeftArrow}{{$\leftarrow$}}1
    {\\LeanLitCompL}{{$\circ_L$}}2
    {\\LeanLitDelta}{{$\delta$}}1
    {\\LeanLitBeta}{{$\beta$}}1
    {\\LeanLitAlpha}{{$\alpha$}}1
    {\\LeanLitLambda}{{$\Lambda$}}1
    {\\LeanLitSubZero}{{$_0$}}1
    {\\LeanLitSubOne}{{$_1$}}1
}
\crefname{lstlisting}{formalization}{formalizations}
\Crefname{lstlisting}{Formalization}{Formalizations}
\crefname{listing}{formalization}{formalizations}
\Crefname{listing}{Formalization}{Formalizations}
\renewcommand{\lstlistingname}{Formalization}

\title{Did We Really Formalize Davis--Kahan? Semantic Alignment for LLM-Assisted Lean Formalization}
\author{%
  Jonathan Crall \quad Brian Hu \\
  Kitware, Inc.
  \And
  Edward Wang \quad Carey E. Priebe \\
  Johns Hopkins University
}
\date{}

\begin{document}
\maketitle

\begin{abstract}
We report a retrospective case study of an LLM-assisted Lean formalization of
29 mathematical results tracked from Davis and Kahan's seminal 1970 paper.  All
Lean results reported as established in this paper are compiler validated and
kernel checked.  The project exposed a separate problem: a checked theorem can
still differ from the mathematical claim being attributed to its source.
Several previously accepted statements were returned to review after source comparison,
and the formalization also produced a checked counterexample to
Davis--Kahan Proposition~4.4 together with a checked repair.  We document the
formalization and source-review practice that emerged during the project and
three sine-theta statement boundaries that illustrate different correspondence
failures.  This is an experience report rather than a controlled evaluation of
the workflow: we did not compare against a baseline, conduct blind or
independent expert adjudication, or estimate a residual semantic-error rate.
The case study therefore supports the narrower conclusion that successful
compilation alone did not establish source correspondence in this project.
\end{abstract}

\section{Introduction}

Lean as a formalization tool checks that a proof term establishes the proposition encoded in a theorem
statement \citep{demoura2021lean4,mathlib2020}.  Establishing that the proposition
expresses a source theorem also requires checking its hypotheses, objects, and
scope.  This agreement is called \emph{semantic alignment}
\citep{lu2025formalalign}.  For a research paper, the intended claim includes
definitions, inherited assumptions, and surrounding mathematical context.
Consequently, trusting an LLM-generated formalization requires more than simply
checking if the Lean compiles: the reviewer must also be able to judge whether
the encoded proposition is semantically aligned with the claim in the source.
Chow illustrates the risk with an agent that replaces ``all doohickeys are
excellent'' by ``all special doohickeys are quasi-excellent''
\citep{chow2026mindreading}.

We encountered this semantic alignment problem while formalizing Davis and Kahan's 1970
spectral-perturbation results \citep{davis1970rotation}.  We wanted reusable
bounds for downstream applications, including the paper's real and complex
Hilbert-space settings.  
Source review revealed several formal statements that
did not match the scope of the corresponding source claims: some were weaker,
while others asserted stronger results than Davis and Kahan had originally proved. Such
strengthening of claims may be mathematically valuable, but they do not by themselves
support confidence that the source theorem has been faithfully represented.
This paper
examines how we organized that review when LLMs could produce more Lean than
we could inspect line by line.

\Cref{fig:workflow} summarizes the human-aligned formalization process that we developed to address these
source-correspondence failures.  The source material remains visible through
context gathering, decomposition, foundation search, formalization, and a
separate source-to-statement review.  \Cref{sec:related} reviews related work
in AI-assisted formalization and semantic alignment; \Cref{sec:dk} defines the
scope and notation of the Davis--Kahan development; \Cref{sec:system} explains
the workflow and its human intervention points; \Cref{sec:cases} presents the
worked semantic-alignment examples and Proposition~4.4 counterexample;
\Cref{sec:practitioner-accounts} gives the public practitioner accounts as
context; and \Cref{sec:limitations} states the limits of the retrospective
evidence.

\begin{figure}[H]
  \centering
  \includegraphics[width=0.98\linewidth]{figures/formalization_workflow.png}
  \caption{Formalization workflow used in this case study.  Inputs include the
  source paper, surrounding mathematical context, and the accumulated Lean
  development.  A target passes through context gathering, decomposition,
  foundation search, agent implementation with compiler-driven revision, and a
  separate skeptical source-to-statement review.  Human supervisors can redirect
  the target, add source context, reject a decomposition or statement, and request
  another review at any return through the loop.  Accepted targets contribute
  checked Lean results, source-correspondence records, and reusable foundations;
  failed semantic review returns the target for revision.  The source comparison
  checks for extra or missing hypotheses and other statement drift rather than
  attempting to prove that the source hypotheses are minimal.}
  \label{fig:workflow}
\end{figure}

Our contributions are: (1) a compiler-validated formalization case study
covering all 29 tracked Davis--Kahan results, with 28 currently recorded as proving
their source claims and Proposition~4.4 formally refuted as printed; (2) a
retrospective record of source-to-statement review, including concrete cases in
which previously accepted checked statements were returned to review and repaired; and
(3) a fully formalized Proposition~4.4 counterexample and a proved replacement
for the surviving $Q$-norm class.  We also release a convenience collection of
public practitioner accounts as contextual material, but do not use that
collection to estimate prevalence or to validate the workflow.

The workflow was not designed prospectively as an experiment.  It emerged over
the course of the formalization, and we did not compare it against compile-only,
human-only, or other review baselines.  Nor did we conduct blind review,
measure inter-rater agreement, or obtain independent Davis--Kahan expert
adjudication for the 29 source correspondences.  Our evidence is therefore
retrospective: it records what the project checked, which accepted
correspondence judgments were later returned to review, and how those particular defects
were repaired.  We make no claim that the workflow is optimal or that the final
correspondence judgments have a measured error rate.  In particular, this paper
does not introduce a validated semantic-alignment benchmark, metric, or general
evaluation framework.


\section{Related Work}
\label{sec:related}

Large-language-model support for formal mathematics has progressed from
benchmark-scale theorem proving and autoformalization toward longer-horizon
mathematical work.  AlphaProof demonstrated olympiad-level formal reasoning
\citep{hubert2026alphaproof}, while recent public reports describe frontier
models producing substantial Lean developments: Anthropic reports a largely autonomous
formalization of Fermat's Last Theorem \citep{anthropic2026flt}, and OpenAI
reports an AI-generated Navier--Stokes argument accompanied by a Lean
formalization \citep{openai2026navierstokes}.  These developments make the
human-facing question studied here increasingly practical: after a formal
artifact compiles, a reader still needs evidence that its statement corresponds
to the mathematical claim being attributed to the source.  Broader surveys of
LLM autoformalization and theorem proving provide context for the rapidly
changing system landscape \citep{weng2025survey}.

Beyond these headline demonstrations, several projects treat AI-assisted
formalization as a workflow and systems problem.  LeanFlow organizes iterative
autoformalization around explicit workflow stages \citep{milikic2026leanflow};
FormalScience studies scalable human-in-the-loop agentic formalization
\citep{meadows2026formalscience}; and Collins et al. characterize initial
human--AI proof-formalization workflows \citep{collins2026humanai}.  Other work
targets longer source documents, including mathematical literature and
textbooks \citep{wang2026m2f,gloeckle2026textbook}.

Several systems directly address source-to-formal correspondence.
FormalAlign evaluates informal/formal agreement \citep{lu2025formalalign};
ShadowBench probes formal statements with auxiliary implications
\citep{han2026shadowbench}; EconCSLib presents source text, Lean statements,
and human correspondence judgments \citep{garg2026econcs}; and Lean
Atlas/Lean Compass expose semantic dependencies while pruning proof-only
dependencies \citep{yanahama2026leanatlas}.  LeanMarathon addresses
adversarial review \citep{zhang2026leanmarathon}, FormaTheoria addresses source
reconciliation \citep{nie2026formatheoria}, and LeanArchitect synchronizes
informal blueprints with Lean declarations \citep{zhu2026leanarchitect}.
LeanDojo \citep{yang2023leandojo}, Pantograph \citep{aniva2024pantograph}, and
LeanInteract \citep{poiroux2025leaninteract} provide programmatic access to
Lean.  As our project grew, we developed narrower tools for inspecting
compiler-derived declarations and maintaining source-to-statement records.

\section{What we formalized}
\label{sec:dk}

Davis and Kahan's Section~2 $\sin\Theta$ theorem compares a trial subspace
with an exact spectral subspace of the perturbed full-space operator.  Their
notation distinguishes the reference operator $A$ from the perturbed operator
$A+H$.  To match the Lean signatures below, we use $A$ for the perturbed
full-space operator: the Lean parameter \texttt{A} therefore corresponds to
Davis--Kahan's $A+H$, and their reference operator $A$ is not a separate
parameter of these signatures.  In population/sample terminology, the Lean
$A$ is the perturbed or sample operator, while Davis--Kahan's unperturbed $A$
is the reference or population operator.

Let $A_0$ be a self-adjoint operator on trial coordinates and let $E_0$ embed
those coordinates isometrically into the ambient Hilbert space. The range of
$E_0$ is the \emph{trial subspace}; in applications it is often an approximate
spectral subspace. For a trial vector $x$, one can either embed it and apply
the ambient operator, giving $AE_0x$, or apply the trial operator first and
then embed, giving $E_0A_0x$. The residual

\begin{equation}
  R=AE_0-E_0A_0
  \label{eq:residual}
\end{equation}

measures the discrepancy between these two operations. If $R=0$, the trial
subspace is invariant under $A$, and $A_0$ describes the action of $A$ on that
subspace. In the unbounded setting this identity is imposed on the domain of
$A_0$, and $R$ is a bounded extension of that residual.

Let $F_0$ be isometric coordinates for the desired exact subspace and $F_1$ for
its orthogonal complement, with $\Lambda_1$ the corresponding complementary
spectral block of $A$. The operator $\Theta_0$ records the principal angles
between the trial subspace and the desired exact subspace. The positive
operator

$$
  \cos\Theta_0
  = \bigl(E_0^*F_0F_0^*E_0\bigr)^{1/2}
$$

has, in finite dimensions, the cosines of those principal angles as its
eigenvalues. Since $E_0$ is an isometry, $E_0^*E_0=I$, and the identity
$\sin^2\Theta_0=I-\cos^2\Theta_0$ gives

$$
  \sin\Theta_0
  = \bigl(E_0^*(I-F_0F_0^*)E_0\bigr)^{1/2}.
$$

Thus $\sin\Theta_0$ is a positive operator on trial coordinates whose
eigenvalues are the sines of the principal angles.

Define

\begin{equation}
  S=(I-F_0F_0^*)E_0
  \label{eq:sin-theta-representative}
\end{equation}

Davis--Kahan show that $S$ is $\sin\Theta_0$ followed by an isometry. The two
operators therefore have the same singular values and

$$
  N(S)=N(\sin\Theta_0)
$$

for every unitarily invariant norm $N$. Formalizations~1 and~2 use $S$ directly.
Formalization~3 states the bound on a positive trial-coordinate operator obtained
as the functional-calculus sine of the directed angle.  Lean proves that this
operator is exactly $|S|$, and the symmetric ideal gauge is invariant under
passing from $S$ to $|S|$.

If the trial block $A_0$ and complementary block $\Lambda_1$ are separated by
$\delta>0$, Davis--Kahan prove, for every unitarily invariant norm for which the
two displayed quantities are defined,
\begin{equation}
  \delta\,N(S)\le N(R),
  \qquad\text{equivalently}\qquad
  \delta\,N(\sin\Theta_0)\le N(R).
  \label{eq:sin-theta}
\end{equation}
The familiar separation places one spectrum in a finite interval and the other
outside its $\delta$-neighborhood.  Davis--Kahan also allow the separating
interval to be half-infinite: one block may be bounded above and the other
bounded below by a value at least $\delta$ away, in either order.  The theorem
covers real and complex separable Hilbert spaces and permits unbounded
self-adjoint operators under the stated domain conditions
\citep{davis1970rotation}.  \Cref{app:formalization1} expands this notation.

Our 29 targets comprise the four unnumbered headline theorems in Davis--Kahan
Section~2 and every named theorem, proposition, lemma, and corollary established
in the paper. \Cref{app:data} specifies the counting convention.  At the state
reported here, all Lean code in the formalization builds, and every one of the
29 result rows is marked \texttt{proved\_in\_build}.  The maintained result
register classifies 28 rows as proving their accepted source claims and
Proposition~4.4 as refuted as transcribed, with a separate proved repair.
There are no pending compiler-validation items among the formal results reported
as established here.
\Cref{app:target-status} gives a per-target index of the final source anchor,
formal disposition, correspondence status, and canonical Lean evidence.

Proof checking and source correspondence are tracked separately, so ``29/29''
records formal coverage and a current project-internal correspondence judgment;
it is not a semantic-completeness certificate or an estimate of residual error.
The $\sin\Theta$ example below illustrates the distinction: a checked theorem
can have the expected conclusion and unbounded operator types while still
covering only a strict subset of the source's spectral-gap cases.

The development required projection and principal-angle geometry,
singular-subspace geometry, unitarily invariant norms and operator ideals,
Borel spectral calculus, unbounded self-adjoint operator theory, real/complex
transfer, direct rotations, and Sylvester equations.  Parts of the spectral and
operator-theoretic infrastructure were copied, adapted, or generalized from the
Spectra formalization \citep{bornemann2026spectra}.  Reusable components developed
here are being prepared for Tau Ceti, an AI-authored Lean library downstream of
Mathlib \citep{tauceti2026}.

\section{Formalization process}
\label{sec:system}

The workflow in \Cref{fig:workflow} can be summarized by five recurring
activities.  Here, ``human researchers'' means the authors directing the
project, while ``agents'' means the ChatGPT and Claude sessions used for
implementation, mathematical discussion, search, and review.  The two roles
were deliberately not collapsed into an undifferentiated ``we.''

\begin{itemize}
\setlength{\itemsep}{0.25em}
\item \textbf{Gather Context.}  Human researchers selected the source
  targets and supplied the relevant PDFs and framing.  Agents transcribed
  source passages and drafted prose distillations of the definitions,
  assumptions, and arguments needed for the target.  Humans could request more
  context or correct the framing before the target moved on.

\item \textbf{Decompose.}  Agents proposed intermediate results and
  dependencies that would support the target; human supervisors redirected the
  decomposition when it no longer matched the intended result.  This serves a
  role similar to the blueprints used in other large Lean formalizations
  \citep{liquidblueprint2023,pfrblueprint2024,zhu2026leanarchitect}.

\item \textbf{Find Foundations.}  Agents searched Mathlib, existing Lean
  developments, and the literature before rebuilding an intermediate result.
  Human supervisors could redirect this search or ask for reuse of an existing
  abstraction.  Spectra, for example, supplied spectral and operator-theoretic
  material that we reused or extended \citep{bornemann2026spectra}.  Missing
  foundations were included as formalization obligations.

\item \textbf{Formalize.}  Agents implemented the resulting tasks in Lean and
  revised definitions and proofs in response to compiler feedback.  Lean
  mechanically checked the resulting proof terms.  Human supervisors set
  priorities, questioned unexpected statements or proof directions, and could
  send a target back to an earlier stage.  ChatGPT and Claude were used
  throughout the project; we do not assign fixed roles to the two systems.

\item \textbf{Skeptical Review.}  A review agent separate from the implementing
  pass compared source passages with theorem statements, checking hypotheses,
  mathematical objects, scalar field, dimension, boundedness, norm class, gap
  assumptions, constants, and directionality.  Real, complex, and
  scalar-generic variants were reviewed together to detect semantic differences
  among sibling statements.  Human supervisors inspected disagreements and
  could request another review or a mathematical explanation before accepting a
  correspondence judgment.  These checks did not constitute independent
  subject-matter adjudication of every source interpretation.
\end{itemize}

The stages loop.  A mismatch can require more source context, a different
decomposition, additional foundations, or a revised formal statement.  Human
intervention is possible at every return edge: the researchers can add context,
change the framing, reject a decomposition or statement, redirect the search,
or request another review.  We also periodically reorganized definitions and
intermediate results so foundations developed for one target could be reused by
later ones.  As the project grew, we added tools for inspecting the accumulated
formalization and maintaining source-to-statement records.

This procedure emerged during the project; it was not a preregistered review
protocol.  Reviews were sometimes performed in fresh model sessions or by a
different model from the implementing pass, but targets were not systematically
double-coded, review order was not blinded, and no inter-rater agreement was
measured.  At a pinned 4 September snapshot, the retained Git history contains
14 transitions in which a result previously marked semantically accepted was
returned to a nonaccepted state, affecting 10 distinct result rows.  These are
review-state reversals, not 14 invalid Lean proofs: they include mathematical
scope or representation defects as well as defects in the audit machinery.
Because no mismatch taxonomy was fixed in advance, we report those transitions
as retrospective evidence rather than an alignment-accuracy statistic.
\Cref{app:review-record} gives the supporting chronology.

We provisionally accepted a target when Lean checked its formal treatment and a
separate source-to-statement review found its hypotheses, mathematical objects,
conclusion, and scope consistent with the source.  When a correspondence used a
non-obvious representation, the required identification also had to be
established.  Davis--Kahan Proposition~4.4 instead has a checked counterexample
to the printed claim and a proved repair.  Later review can still expose a
mismatch missed by an earlier pass, so this is an operational stopping rule, not
a certificate of source correspondence.  \Cref{app:system} gives additional
procedure details.

\section{Semantic-alignment case studies}
\label{sec:cases}

The examples below show why the project treated compilation and source
correspondence as separate checks.  The first compares three checked
Davis--Kahan Section~2 sine-theta statement boundaries; the second records a
source proposition for which formalization produced a counterexample and a
repaired theorem.

% Historical provenance for editors (not rendered): the semantic-review packet
% at commit 489c01c2cc992a20d38415b5f827cdc046fe7236 listed
% sinTheta_unbounded_intervalExterior_characterizedWitness_rclike under
% "Canonical Lean declarations" and marked the source-correspondence rows
% claimed_exact. Commit cb3b330b later recorded explicitly that the row had
% named this declaration as the exact source match and that this was wrong
% because the theorem covered only the finite interval/exterior branch while
% Davis--Kahan permit half-infinite separation. The same correction also
% discussed an older ideal-membership interpretation; do not use that point as
% the worked mismatch because the project's norm-boundary interpretation later
% changed. The durable defect used here is gap scope.
\subsection{Three sine-theta theorem statements}
\label{sec:scope-mismatch}

The development contains several checked declarations with the factor-one
$\sin\Theta$ inequality in \cref{eq:sin-theta}.  Davis--Kahan first state the
theorem with a finite separating interval and later extend the same estimate to
half-infinite intervals in the Appendix to Section~6.  The three signatures
below differ in how closely their hypotheses follow that complete proved scope.

\paragraph{Formalization 1.}
Formalization~1 is weaker than the target statement because it restricts the
gap to the finite interval/exterior case.  \Cref{lst:sintheta-mismatch}
requires finite $\beta\le\alpha$ and one spectrum to lie in $[\beta,\alpha]$
while the other lies outside $(\beta-\delta,\alpha+\delta)$, with the two roles
interchangeable.

In the Appendix to Section~6, Davis--Kahan extend the same sine-theta estimate
to half-infinite separating intervals.  In those cases one spectral block may be
only semibounded above and the other only semibounded below, with a gap of at
least $\delta$ between them.  Such an instance need not put either block inside
any finite interval $[\beta,\alpha]$.  Formalization~1 therefore covers the
initial finite-interval statement but not the later proved extension.

% Presentation note (not rendered): scripts/build_evidence.py extracts the
% historical signature from commit 489c01c2cc992a20d38415b5f827cdc046fe7236,
% writes the exact text to generated/historical_sin_theta_gap_mismatch_exact.lean,
% and verifies the contemporaneous/corrective semantic-review text from Git.
% The listings input differs only by LeanLit... ASCII typography sentinels.
\begin{nolinenumbers}
\begin{minipage}{\linewidth}
\lstinputlisting[
  style=leanpaper,
  firstline=6,
  caption={The finite interval/exterior form matches the initial Section~2 gap condition but does not include the half-infinite extension proved later.  The hypothesis \texttt{hSinTheta0} names the rectangular representative $S$ in \cref{eq:sin-theta-representative}.},
  label={lst:sintheta-mismatch}
]{generated/historical_sin_theta_gap_mismatch_presentation.lean}
\end{minipage}
\end{nolinenumbers}

\paragraph{Formalization 2.}
Formalization~2 is stronger than the target statement on the operator
hypotheses because those hypotheses are weaker.  \Cref{lst:stronger-sintheta}
has no ambient separability assumption.  Its
\texttt{IsTrialResidual} predicate requires only that $E_0$ carry
$\operatorname{dom}A_0$ into the domain of the full-space operator, rather than
requiring equality of the two pulled-back domains.  Its ordered
\texttt{FormBoundedSylvesterGap} branches use quadratic-form semibounds, which
are implied by the source's spectral half-line conditions.  On these axes the
theorem applies to every source instance and to additional instances satisfying
the weaker conditions.

Its norm interface is expressed differently from the source convention:
\texttt{hR : N.Mem R} assumes that the residual has finite $N$-gauge, and the
conclusion also proves membership of the sine representative in that gauge's
domain.  The greater generality comes from the weaker operator hypotheses;
the norm interface is a separate difference from the source convention.

% Editor provenance (not rendered): this display is generated verbatim from
% sinTheta_unbounded_formGap_symmetricNorming_rclike in
% DavisKahan/Sources/DavisKahan1970/SineTheta/Presentation.lean.  The exact
% Unicode signature is generated/stronger_sin_theta_exact.lean.
\begin{nolinenumbers}
\begin{minipage}{\linewidth}
\lstinputlisting[
  style=leanpaper,
  firstline=6,
  caption={This theorem uses weaker operator-side hypotheses: ambient separability is absent, the source common-domain equality is weakened to forward domain inclusion, and the ordered gap cases use form semibounds.  The angle term in \texttt{N.gauge} is the same rectangular representative $S$ named by \texttt{hSinTheta0} in Formalization~1.},
  label={lst:stronger-sintheta}
]{generated/stronger_sin_theta_presentation.lean}
\end{minipage}
\end{nolinenumbers}

\paragraph{Formalization 3.}
Formalization~3 is the closest of the three to the target statement.
\Cref{lst:current-sintheta} restores the source's separable ambient scope and
states the left-hand side on \texttt{sourceDirectedSinThetaOperator E0 F0}, the
positive trial-coordinate sine operator.  It also models the source convention
that a displayed unitarily invariant norm statement is vacuous when either norm
is undefined.  The two \texttt{N.Mem} implications after the colon restrict the
numerical comparison to the cases in which both displayed norms exist.

% Editor provenance (not rendered): this display is generated verbatim from
% sinTheta_unbounded_formGap_whereDefinedUIN_rclike in
% DavisKahan/Sources/DavisKahan1970/SineTheta/Presentation.lean. The exact
% Unicode signature is generated/current_sin_theta_exact.lean; the listings
% input differs only by ASCII typography sentinels. Source fidelity is a ledger
% judgment, not a theorem-name claim.
\begin{nolinenumbers}
\begin{minipage}{\linewidth}
\lstinputlisting[
  style=leanpaper,
  firstline=6,
  caption={The statement closest to Davis--Kahan restores ambient separability, states the norm comparison where both gauges are defined, and places the left-hand side on the positive trial-coordinate $\sin\Theta_0$ operator.  The implementation identifies that operator with $|S|$ for $S=(I-F_0F_0^*)E_0$.},
  label={lst:current-sintheta}
]{generated/current_sin_theta_presentation.lean}
\end{minipage}
\end{nolinenumbers}

Formalization~3 retains the forward-domain and form-gap hypotheses of
Formalization~2.  After translating Davis--Kahan's $A+H$ to the Lean parameter
$A$, their unbounded discussion states the common-domain condition
\[
  \{x:E_0x\in\operatorname{dom}A\}=\operatorname{dom}A_0.
\]
This equality implies the forward inclusion used by \texttt{IsTrialResidual};
the source's spectral half-line assumptions likewise imply the ordered form
semibounds.  These are established parts of the compiled development, not
pending replacements: the displayed Formalization~3 declaration and the bridge
from the source common-domain condition to its Lean trial-residual hypothesis
are both compiler validated.  \Cref{app:formalization1,app:formalization2,app:formalization3}
gives a reader's guide to the three signatures and the sine-angle norm
representation.

\subsection{A counterexample to Davis--Kahan Proposition 4.4}
\label{sec:prop44}

An additional outcome of this effort was the discovery that
Davis--Kahan Proposition~4.4 is false as printed. It asserts that, for real
Hilbert spaces with largest principal angle at most $\pi/3$, the direct
rotation minimizes $\lVert I-V\rVert$ for every unitarily invariant norm
\citep{davis1970rotation}. We obtain an explicit counterexample in
$\mathbb{R}^4$ for the trace norm. With two principal angles equal to $\pi/4$,
the direct rotation has trace displacement
$4\sqrt{2-\sqrt{2}}$, while an admissible orthogonal competitor $W$ has
displacement $2\sqrt{2}$.

The same witness identifies the failing step in the printed proof. Write
$K=I-W$, and let $\Omega_1$ and $\Omega_2$ project onto the two orthogonal
principal planes. Davis--Kahan equation~(4.3) asserts, in this case,

$$
  \lVert K\rVert_4
  \geq
  \lVert K\Omega_1\rVert_2+\lVert K\Omega_2\rVert_2,
$$

where the subscripts denote Ky Fan norms: $\lVert A\rVert_j$ is the sum of
the $j$ largest singular values of $A$. The inference treats the orthogonal
decomposition of the domain into principal planes as though the corresponding
singular-value sums could be added. Orthogonality of
$\Omega_1E$ and $\Omega_2E$, however, does not imply orthogonality of
$K\Omega_1E$ and $K\Omega_2E$. When these image spaces overlap, the singular
values of the two restricted operators need not combine to give those of
$K$.

That is exactly what happens in the counterexample:

$$
  \lVert K\rVert_4=2\sqrt{2},
  \qquad
  \lVert K\Omega_1\rVert_2
  =
  \lVert K\Omega_2\rVert_2
  =2.
$$

Thus Davis--Kahan equation~(4.3) would require $2\sqrt{2}\geq4$. The subsequent blockwise
estimate, where the $\pi/3$ hypothesis enters, holds on this witness; the
failure occurs earlier, when the separate block estimates are combined.
Moreover, the construction extends to every $0<\theta<\pi/2$, so no smaller
positive angle bound repairs the proposition.

The repair is instead to restrict the norm class. A unitarily invariant norm
$N$ is a \emph{$Q$-norm} if
$N(A)^2=M(A^*A)$ for some unitarily invariant norm $M$
\citep{bolotnikov2001minimal}. Applying the valid Davis--Kahan
Proposition~4.3 to $M$ and taking square roots shows that, for every acute pair
of subspaces, the direct rotation minimizes $\lVert I-V\rVert$ for every
$Q$-norm. This holds in finite-dimensional real or complex inner-product
spaces and includes the Schatten $p$-norms for $p\geq2$. Thus the repair
retains acuteness, removes the $\pi/3$ and real-space restrictions, and narrows
the class of norms.  The complete chain is formalized in Lean: the printed
Proposition~4.4 is represented as a proposition, an explicit $\mathbb R^4$
witness proves its negation, and the $Q$-norm replacement is a separate checked
theorem.  \Cref{app:target-status} identifies the corresponding declarations.

To the best of our knowledge, Davis--Kahan Proposition~4.4 has not previously
been explicitly identified as false under its stated hypotheses. There were
close precedents. Bolotnikov, Li, and Rodman showed that the corresponding
minimization can fail for Schatten $p<2$ norms, including the trace norm, but
their explicit obstruction occurs at principal angle $\pi/2$, outside the
Davis--Kahan short-angle regime \citep{bolotnikov2001minimal}. Qiu, Zhang, and
Li later repeat the Davis--Kahan $\pi/3$ minimality claim and, elsewhere in the
same paper, compare the angle vectors $(\pi/4,\pi/4)$ and $(\pi/2,0)$ under an
additive sine functional, exhibiting the same concentration mechanism without
connecting it to Davis--Kahan Proposition~4.4 \citep{qiu2006grassmann}. Our
literature search found no earlier correction or counterexample within the
stated real $\theta_{\max}\leq\pi/3$ regime.


\section{Context: public practitioner accounts}
\label{sec:practitioner-accounts}

We release a convenience collection of public accounts of AI-assisted Lean
work because these reports are otherwise scattered across papers, blogs,
repositories, forums, and social posts.  We do not treat the collection as a
representative survey, an evaluation dataset, or evidence that the workflow in
\cref{sec:system} is typical or effective.  It was assembled through
LLM-assisted web search, and the screening log was introduced only after part
of the collection already existed.

An account is included when a public first-person source describes an
AI-assisted Lean project or workflow in enough detail to support source-grounded
coding of at least some supervision, review, or semantic-checking fields.
Status-only reports without sufficient workflow detail are held rather than
coded; structured human--AI studies and adjacent system papers are tracked
separately.  Multiple sources can support one project/workflow episode, and
known repeated practitioners are clustered rather than being treated as
independent people.  The current snapshot contains \PractitionerAccountCount{}
practitioner accounts supported by \PractitionerSourceCount{} source records;
\Cref{app:data} documents the schemas, screening dispositions, and known
selection limitations.

The collection is useful primarily as a directory of concrete practice.  It
contains examples ranging from close reading of generated definitions and
statements to workflows in which humans inspect little Lean directly, as well
as reports where separate review agents or human counterexamples exposed
specification problems
\citep{kahle2026keeptrying,yadav2026verification,gowers2026leiden,zhang2026slt,zaliva2026semantics}.
Those examples motivated questions we asked of our own process, but their
frequency in this convenience sample should not be interpreted as prevalence
among Lean users.

\section{Limitations}
\label{sec:limitations}

This paper is a retrospective experience report about one formalization project
in one area of operator theory.  The workflow emerged while the project was in
progress; we did not prospectively define treatments, baselines, or success
metrics, and therefore cannot conclude that it is more reliable than a simpler
``LLM + compile + manual spot check'' process, human-only review, or another
agentic workflow.  The retained acceptance reversals show that compilation and
prior internal acceptance were insufficient in this project, not that the
particular review process is optimal.

The source-correspondence judgments are also not independently calibrated.  The
human supervisors were not Davis--Kahan specialists, targets were not reviewed
blindly by multiple independent raters, no inter-rater agreement was measured,
and we did not obtain external subject-matter-expert adjudication for a fixed
subset.  Fresh model sessions and repeated hostile reviews supplied useful
checks, but they do not substitute for those forms of validation.  We therefore
report current correspondence status and review history without an alignment
accuracy or residual-error estimate.

The 29-result denominator is a defined register of named and headline results,
not a proof that every semantic aspect of the 1970 paper has been captured.
Later review has repeatedly returned previously accepted rows to review.  To mitigate
that risk for downstream users, the artifact keeps compiler status, formal
disposition, and source-correspondence status separate; records the source
boundary used for each result; and requires explicit bridges for non-obvious
representation changes.  These practices make later correction auditable, but
they do not make the current semantic judgments infallible.

Finally, the practitioner-account collection is opportunistic and
nonrepresentative.  Its value is archival and contextual: it centralizes public
first-person reports and retains enough provenance for readers to inspect the
coding.  It should not support population-level claims about AI-assisted
formalization practice.

\section{Conclusion}

\paragraph{Did we formalize Davis--Kahan?}
All Lean code reported here compiles, and the maintained register has a checked
formal treatment for each of the 29 tracked results: 28 proved results and one
formal refutation, Proposition~4.4, with a separate proved repair.  The three
sine-theta signatures separate mathematical strength from correspondence with
a particular source statement.  One is a strict specialization because its gap
hypothesis is too narrow; another has weaker operator-side hypotheses than
Davis--Kahan; and the closest current boundary restores the source's separable
scope and where-defined norm convention while retaining two operator-side
generalizations.  What remains unmeasured is semantic residual risk: the
project has no calibrated stopping rule or independent expert audit establishing
that every current source-to-statement judgment is correct.

The case study therefore supports reviewing each conclusion together with the
hypotheses and definitions that determine its scope.  A matching inequality,
appropriate operator types, and a successful Lean proof do not by themselves
establish that the encoded proposition covers the same mathematical cases as
the source.

\paragraph{AI assistance.}
ChatGPT and Claude were used throughout the formalization and manuscript
preparation, including code generation, mathematical discussion, literature
search, and source review. Retained token telemetry and derived resource estimates are
reported in \Cref{app:resources}.

\ifnum\PublicVersion=1\relax
\section*{Acknowledgements}
This material is based upon work supported by the Defense Advanced Research
Projects Agency (DARPA) under Contract No. HR001125CE017. Any opinions, findings,
and conclusions or recommendations expressed in this material are those of the
author(s) and do not necessarily reflect the views of the Defense Advanced
Research Projects Agency (DARPA).
\fi

\clearpage
\begingroup
\sloppy
\bibliographystyle{unsrtnat}
\bibliography{references}
\endgroup

\clearpage
\appendix
\section{Supplementary data}
\label{app:data}

The supplementary tables contain \PractitionerAccountCount{} public
practitioner accounts---public posts, papers, discussions, or project notes in
which practitioners describe their own use of AI in Lean---supported by
\PractitionerSourceCount{} practitioner-source records.  Known repeated practitioners form
\PractitionerOverlapClusterCount{} overlap clusters; one controlled human--AI
workflow study is kept separate.  Accounts are project/workflow episodes rather
than unique-person records, and multiple sources may support one account.  The
collection was assembled through LLM-assisted web search and is not a systematic
or representative sample.

The screening table records \PractitionerScreenedCandidateCount{} candidates and
their dispositions.  It begins on 8 September 2026 and does not reconstruct
every source encountered in earlier searches.

The supplementary material also contains the 29-result register used for the
Davis--Kahan formalization.  It counts the four unnumbered headline theorems in
Davis--Kahan Section~2 and every named theorem, proposition, lemma, and corollary
established in the source.  Definitions, proof steps, examples, immediate
consequences, and open or deferred claims are not counted separately.
Davis--Kahan Proposition~4.4 remains a counted result because its formal
treatment gives a counterexample to the printed claim and proves a repaired
theorem.

\ifshowartifacturl
Project materials: \href{\artifacturl}{public GitHub repository}.
\else
Identifying repository URLs are suppressed in this double-blind build and are
included in the supplementary artifact.
\fi

\section{Formalization procedure}
\label{app:system}

\Cref{sec:system} summarizes the workflow.  Source transcriptions, prose
summaries, result records, and review judgments were kept outside individual
model conversations so work could continue across sessions.  The formalization
used a pinned Lean/Mathlib environment and reusable local foundations; exact
dependency revisions and helper commands are supplied with the artifact rather
than repeated here.  Mechanical checks established the Lean declarations and
their types, while correspondence judgments were recorded separately.  Source
review included inherited section assumptions and the fields of structures
occurring in a theorem type, because both can restrict a claim without changing
its concluding inequality.

Source review compared each tracked Davis--Kahan result with the elaborated Lean
statement and definitions that affected its meaning.  Recurring questions were
the scalar field, dimension, bounded or unbounded operator scope, norm class,
gap conditions, constants, directionality, and the identity of angle and
residual operators.  Mechanical checks established properties of the formal
artifact; they did not determine whether our reading of the source was correct.

ChatGPT and Claude were used for implementation, mathematical discussion,
search, and review.  Human supervision consisted mainly of monitoring progress,
supplying framing, asking models to explain inconsistencies, and requesting
additional or independent reviews.  The humans directing the formalization were
not Davis--Kahan subject-matter experts, so disagreements about source meaning
were often investigated by probing the models in different ways rather than by
an independent expert reading of the theorem.

\section{Davis--Kahan sine-theta notation and Formalization 1}
\label{app:formalization1}

The sine-theta theorem is naturally expressed with coordinate maps.  Davis and
Kahan write the reference operator as $A$ and the perturbed full-space operator
as $A+H$.  The Lean signatures use the parameter name \texttt{A} for the latter:
throughout this reader's guide, $A$ therefore denotes Davis--Kahan's $A+H$.
Their unperturbed reference operator does not appear as a separate parameter in
these signatures.

The trial-coordinate space carries a possibly unbounded self-adjoint operator
$A_0$, and the isometry $E_0$ embeds those coordinates into the ambient Hilbert
space.  Its range is the trial subspace.  The bounded residual $R$ extends
\[
  AE_0-E_0A_0
\]
from the domain of $A_0$.  The desired exact subspace is represented by an
isometry $F_0$; $F_1$ represents its orthogonal complement, and the
self-adjoint operator $\Lambda_1$ is the complementary spectral block.  The
exact-coordinate hypotheses say that the ranges of $F_0$ and $F_1$ are
orthogonal and complete and that $F_1$ intertwines $\Lambda_1$ with $A$.

The source's positive operator $\sin\Theta_0$ acts on the trial-coordinate
space.  The rectangular map
\[
  S=(I-F_0F_0^*)E_0
\]
sends trial coordinates into the ambient space and extracts the component
orthogonal to the desired exact subspace.  Davis--Kahan explicitly observe that
their corresponding rectangular operator is $\sin\Theta_0$ followed by an
isometry.  It therefore has the same singular values
$\sin\theta_1,\sin\theta_2,\ldots$ and the same value under every unitarily
invariant norm:
\[
  N(S)=N(\sin\Theta_0).
\]
Formalizations~1 and~2 express the angle term through $S$.  Formalization~3 uses
a positive trial-coordinate operator defined as the functional-calculus sine of
the directed angle and proves that this operator equals $|S|$.

Formalization~1 makes this representation visible by introducing a variable
named \texttt{sinTheta0} and the hypothesis \texttt{hSinTheta0}.  The hypothesis
defines that variable to be the rectangular map $S$.  The source's positive
$\sin\Theta_0$ and this Lean map have different codomains; equality of their
unitarily invariant norms follows from their common singular values.  This is
the identification needed to read \texttt{N.gauge sinTheta0} as the source term
$N(\sin\Theta_0)$.

\paragraph{Lean operator notation.}
The signatures use three operator notations repeatedly.  For normed spaces
$E$ and $F$ over the scalar field $\mathbb K$,
\[
  E\mathrel{\to}\!\mathtt{L}[\mathbb K]F
\]
denotes a continuous linear map, hence an everywhere-defined bounded linear
operator.  The notation
\[
  E\mathrel{\to}\!\mathtt{L}_{\mathrm p}[\mathbb K]F
\]
denotes a \texttt{LinearPMap}: a linear operator together with an explicitly
tracked domain.  This is the type used for the possibly unbounded operators
$A$, $A_0$, and $\Lambda_1$.  Thus the $\mathtt L$ and $\mathtt L_{\mathrm p}$
forms distinguish bounded everywhere-defined maps from domain-aware partial
maps.

For continuous linear maps, $S\mathbin{\circ_L}U$ denotes composition,
$x\mapsto S(Ux)$.  In particular,
\[
  (I-F_0\mathbin{\circ_L}F_0^*)\mathbin{\circ_L}E_0
\]
is the Lean expression for the rectangular sine representative
$(I-F_0F_0^*)E_0$.  The suffix $L$ on $\circ_L$ identifies composition in the
continuous-linear-map API.  The notation \texttt{F0.adjoint} denotes the
adjoint $F_0^*$.

In Lean identifiers, \texttt{Mem} abbreviates \emph{membership}. Thus
\texttt{N.Mem T} means that the operator $T$ belongs to the operator ideal, or
equivalently to the domain on which $N$ has a finite gauge.

Throughout these signatures, $\mathbb K$ is the scalar type and
\texttt{RCLike} supplies the real-or-complex scalar structure used by the
generic theorem.

Formalization~1 is scalar-generic over real and complex Hilbert spaces and
uses potentially unbounded operators.  The two alternatives in its
\texttt{hspectral} hypothesis say that one spectrum lies in a finite interval
$[\beta,\alpha]$ and the other lies
outside the enlarged open interval
$(\beta-\delta,\alpha+\delta)$.  Swapping the two blocks gives the second
alternative.  Davis--Kahan also allow the separating interval to be
half-infinite, so neither block need have spectrum contained in a finite
interval.  Those cases cannot in general satisfy Formalization~1.

The historical signature also contains two scalar-generic capability classes.
They package analytic results used by the proof and have instances for the real
and complex source fields.  Its norm API also predates the current where-defined
interface.  The finite interval/exterior hypothesis is the restriction relevant
to the source-scope comparison here.

\begin{table}[H]
\centering
\small
\caption{Reader's guide to Formalization~1.  This table establishes the
baseline notation used again in Formalizations~2 and~3.}
\label{tab:formalization1-glossary}
\begin{tabularx}{\linewidth}{@{}>{\raggedright\arraybackslash}p{0.41\linewidth}>{\raggedright\arraybackslash}X@{}}
\toprule
Lean expression & Mathematical role \\
\midrule
\texttt{[RCLike }$\mathbb K$\texttt{]}
  & The scalar field may be real or complex. \\
$\texttt{A}:E\mathrel{\to}\!\mathtt{L}_{\mathrm p}[\mathbb K]E$,
$\texttt{A0}:F\mathrel{\to}\!\mathtt{L}_{\mathrm p}[\mathbb K]F$,
$\Lambda_1:G\mathrel{\to}\!\mathtt{L}_{\mathrm p}[\mathbb K]G$
  & The perturbed full-space operator $A$ (Davis--Kahan's $A+H$), trial
    operator $A_0$, and complementary exact self-adjoint block $\Lambda_1$;
    each may be unbounded. \\
$\texttt{E0}:F\mathrel{\to}\!\mathtt{L}[\mathbb K]E$,
$\texttt{F0}:H\mathrel{\to}\!\mathtt{L}[\mathbb K]E$,
$\texttt{F1}:G\mathrel{\to}\!\mathtt{L}[\mathbb K]E$
  & Bounded coordinate maps for the trial subspace, desired exact subspace,
    and its orthogonal complement. \\
$\texttt{R}:F\mathrel{\to}\!\mathtt{L}[\mathbb K]E$
  & The bounded residual extending $AE_0-E_0A_0$ from $\operatorname{dom}A_0$. \\
\texttt{IsTrialResidual A A0 E0 R}
  & $E_0$ is isometric, maps $\operatorname{dom}A_0$ into
    $\operatorname{dom}A$, and $R=AE_0-E_0A_0$ on that domain. \\
\texttt{IsExactSpectralDecomposition A }$\Lambda_1$\texttt{ F0 F1}
  & $F_0,F_1$ are orthogonal isometric coordinates forming a complete
    decomposition, and $F_1\Lambda_1=AF_1$ on $\operatorname{dom}\Lambda_1$. \\
\texttt{hSinTheta0}
  & Defines the displayed rectangular sine representative as
    $S=(I-F_0F_0^*)E_0$; $S$ has the same singular values and unitarily
    invariant norms as the source's positive $\sin\Theta_0$. \\
$\beta,\alpha,\delta$, \texttt{hspectral}
  & The finite interval/exterior gap; it does not express the two
    half-infinite source configurations. \\
\texttt{hR : N.Mem R}
  & The residual is in the domain of the gauge used in the conclusion. \\
\bottomrule
\end{tabularx}
\end{table}

\section{Reading Formalization 2}
\label{app:formalization2}

Formalization~2 retains Formalization~1's scalar-generic spaces, operator
types, residual predicate, and exact decomposition.  It replaces the finite
\texttt{hspectral} condition by \texttt{FormBoundedSylvesterGap}, which includes
the finite interval/exterior case and both ordered semibounded cases.  In the
ordered cases, the source's spectral half-line hypotheses imply the form
semibounds used in Lean.  As in Formalization~1, no ambient separability
assumption is imposed and \texttt{IsTrialResidual} uses forward domain inclusion
rather than the source common-domain equality.

The norm parameter is a \texttt{SymmetricNormingFunction}.  The theorem assumes
\texttt{hR : N.Mem R} and concludes both
\[
  N.\mathrm{Mem}(S)
  \qquad\text{and}\qquad
  \delta N(S)\le N(R).
\]
The first conclusion transfers membership in the norm ideal from the residual
to the sine representative.  Davis--Kahan state the numerical comparison under
their convention that a displayed norm result is vacuous when a norm does not
exist.  This Lean norm interface instead adds a residual-membership premise and
a sine-representative membership conclusion.  Relative to Davis--Kahan,
Formalization~2 weakens the operator-side assumptions; the norm interface
changes separately.

The expression inside \texttt{N.gauge} is exactly the map $S$ defined above.
Formalization~1 exposes the same map through \texttt{hSinTheta0};
Formalization~2 inlines it.  Since $S$ and the source's positive
$\sin\Theta_0$ have the same singular values, \texttt{N.gauge S} is the norm
appearing in the source inequality.

\Cref{tab:formalization2-glossary} translates the signature.

\begin{table}[H]
\centering
\small
\caption{Changes introduced by Formalization~2 relative to the baseline in
\cref{tab:formalization1-glossary}.  Unchanged operator and coordinate-map
notation is omitted.}
\label{tab:formalization2-glossary}
\begin{tabularx}{\linewidth}{@{}>{\raggedright\arraybackslash}p{0.42\linewidth}>{\raggedright\arraybackslash}X@{}}
\toprule
Lean expression & Mathematical role \\
\midrule
\texttt{N : SymmetricNormingFunction}
  & A symmetric norming function and its associated operator-ideal gauge. \\
\texttt{hgap : FormBoundedSylvesterGap }$A_0$ $\Lambda_1$ $\delta$
  & Replaces the finite interval/exterior hypothesis by the full gap family:
    the finite case or either ordered half-infinite configuration. \\
\texttt{N.Mem }$S$ $\land$ $\delta$\texttt{ * N.gauge }$S$
$\le$ \texttt{N.gauge R}
  & In addition to $\delta N(S)\le N(R)$, the conclusion places $S$ in the
    gauge domain when \texttt{hR : N.Mem R} holds. \\
\bottomrule
\end{tabularx}
\end{table}

\section{Reading Formalization 3}
\label{app:formalization3}

Formalization~3 keeps the operator and gap predicates of Formalization~2,
adds the separable ambient hypothesis stated by Davis--Kahan, and uses a norm
object with an explicit domain of definition.  Its conclusion is written as two
implications,
\[
  N(\sin\Theta_0)\ \text{defined},\qquad N(R)\ \text{defined}
  \quad\Longrightarrow\quad
  \delta N(\sin\Theta_0)\le N(R).
\]
The Lean term \texttt{sourceDirectedSinThetaOperator E0 F0} is the positive
trial-coordinate sine operator.  It is defined by applying sine through
continuous functional calculus to the directed angle reconstructed on
$[0,\pi/2]$.  Lean proves that it equals $|S|$ for
$S=(I-F_0F_0^*)E_0$.  Polar decomposition then gives the same symmetric ideal
membership and gauge for $S$ and $|S|$.

\begin{table}[H]
\centering
\small
\caption{Changes introduced by Formalization~3 relative to Formalization~2.
The operator types, coordinate maps, residual predicate, exact decomposition,
and form-gap predicate are unchanged.}
\label{tab:formalization3-glossary}
\begin{tabularx}{\linewidth}{@{}>{\raggedright\arraybackslash}p{0.42\linewidth}>{\raggedright\arraybackslash}X@{}}
\toprule
Lean expression & Mathematical role \\
\midrule
\texttt{[SeparableSpace E]}
  & Restricts the perturbed full-space Hilbert space to the separable setting
    used in the source. \\
\texttt{N : NormalizedSymmetric}\newline
\texttt{OperatorIdealFamily}
  & A normalized unitarily invariant norm object together with the set of
    operators on which its gauge is defined. \\
\texttt{N.Mem }$S$ $\to$ \texttt{N.Mem R} $\to$\newline
$\delta$\texttt{ * N.gaugeReal }$S$ $\le$ \texttt{N.gaugeReal R}
  & States the numerical comparison when both displayed gauges are defined;
    \texttt{N.gaugeReal} is the real-valued gauge, so the left side is
    $\delta N(\sin\Theta_0)$. \\
\bottomrule
\end{tabularx}
\end{table}

\paragraph{The common-domain condition.}
After translating Davis--Kahan's $A+H$ to the Lean parameter $A$, their
unbounded discussion states the trial-domain condition as
\[
  \{x:E_0x\in\operatorname{dom}A\}=\operatorname{dom}A_0.
\]
Formalization~3's \texttt{IsTrialResidual} predicate uses only the forward
inclusion together with the residual identity on $\operatorname{dom}A_0$.
The repository's \texttt{HasCommonDomain} predicate records the source equality
and proves that it implies this forward inclusion.  Consequently the displayed
Lean theorem is more general on this axis: every source instance satisfies its
domain hypothesis, and the theorem also applies to inputs satisfying only the
forward inclusion.

\paragraph{Norm representation.}
The norm record contains a symmetric ideal gauge, rank-one normalization, and a
where-defined Fan comparison.  Davis--Kahan use the corresponding Ky Fan
comparison principle for unitarily invariant norms.  The formal theorem packages
that comparison as part of the norm interface rather than deriving it from a
smaller collection of norm axioms inside this theorem.

\section{Lean publication activity}

\begin{figure}[H]
  \centering
  \input{generated/lean_activity_timeline_tikz.tex}
  \caption{Monthly publication counts reproduced from the Papers With Lean
  public statistics series, January 2025 through August 2026
  \citep{paperswithlean2026stats}.  Counts are grouped by the captured data's
  \texttt{published} month, and only complete months are shown.  The fraction
  that used AI is unknown.}
  \label{fig:lean-activity}
\end{figure}



\section{Semantic-review page}
\label{app:review-page}

\IfFileExists{figures/semantic-alignment-sine-theta-row.png}{
\begin{figure}[H]
  \centering
  \includegraphics[width=\linewidth,trim=0 250 0 0,clip]{figures/semantic-alignment-sine-theta-row.png}
  \caption{Example semantic-review interface.  The source passage and inherited
  context appear at left, the claimed source-to-Lean correspondence in the
  middle, and compiler-derived Lean evidence at right.}
  \label{fig:dashboard}
\end{figure}
}{}

The interface consolidates source passages, correspondence judgments, and
compiler-derived signatures previously inspected through files and Lean
queries.  Understanding a signature can also require the definitions of its
structures, predicates, coercions, and typeclass assumptions.  Lean Compass's
semantic dependency view suggests how these definitions could be exposed while
omitting proof-only dependencies \citep{yanahama2026leanatlas}.

\section{Resource utilization}
\label{app:resources}

To estimate the cost of the formalization effort, we maintained a resource
accounting record.  Historical installation, attribution, and backfill
limitations make the retained telemetry incomplete.  Explicit backfills are
reported separately and excluded from commit-local calibration and
extrapolation.  The serving-energy row is an estimate derived from the retained
observations.  The energy estimate excludes training, embodied
hardware, client devices, and network energy.

\begin{table}[H]
\centering
\scriptsize
\caption{Retained observed telemetry and derived estimates.  The telemetry is incomplete; the energy value is a modeled order-of-magnitude estimate.}
\begin{tabularx}{\linewidth}{@{}p{0.28\linewidth}p{0.17\linewidth}>{\raggedright\arraybackslash}X@{}}
\toprule
Quantity & Retained / estimate & Basis / interpretation \\
\midrule
\input{generated/resource_snapshot_table.tex}
\end{tabularx}
\end{table}


\begin{table}[H]
\centering
\scriptsize
\caption{AI systems and model families used during the formalization.}
\label{tab:formalization-systems}
\begin{tabularx}{\linewidth}{@{}p{0.24\linewidth}>{\raggedright\arraybackslash}X@{}}
\toprule
System & Model families observed \\
\midrule
\input{generated/model_systems_table.tex}
\end{tabularx}
\end{table}

\begin{table}[H]
\centering
\scriptsize
\caption{Recovered model-resolved token measurements.  These are retained measurements for the model labels with
recovered telemetry; models in \cref{tab:formalization-systems} that do not
appear here were still used.}
\label{tab:model-tokens-appendix}
\resizebox{\linewidth}{!}{\input{generated/model_token_table.tex}}
\end{table}



\clearpage
\section*{NeurIPS Paper Checklist}




\begin{enumerate}

\item {\bf Claims}
    \item[] Question: Do the main claims made in the abstract and introduction accurately reflect the paper's contributions and scope?
    \item[] Answer: \answerYes{}
    \item[] Justification: The abstract and introduction state the case-study scope and distinguish a checked Lean result from source correspondence. The main text and conclusion also state the limits of the current semantic-review evidence.
    \item[] Guidelines:
    \begin{itemize}
        \item The answer \answerNA{} means that the abstract and introduction do not include the claims made in the paper.
        \item The abstract and/or introduction should clearly state the claims made, including the contributions made in the paper and important assumptions and limitations. A \answerNo{} or \answerNA{} answer to this question will not be perceived well by the reviewers. 
        \item The claims made should match theoretical and experimental results, and reflect how much the results can be expected to generalize to other settings. 
        \item It is fine to include aspirational goals as motivation as long as it is clear that these goals are not attained by the paper. 
    \end{itemize}

\item {\bf Limitations}
    \item[] Question: Does the paper discuss the limitations of the work performed by the authors?
    \item[] Answer: \answerYes{}
    \item[] Justification: The paper explicitly states that the workflow is a retrospective single-project case study without a controlled baseline, blind review, inter-rater agreement, independent Davis--Kahan expert adjudication, or a calibrated residual-error estimate.  It also treats the practitioner-account collection as an opportunistic convenience sample rather than a representative survey.
    \item[] Guidelines:
    \begin{itemize}
        \item The answer \answerNA{} means that the paper has no limitation while the answer \answerNo{} means that the paper has limitations, but those are not discussed in the paper. 
        \item The authors are encouraged to create a separate ``Limitations'' section in their paper.
        \item The paper should point out any strong assumptions and how robust the results are to violations of these assumptions (e.g., independence assumptions, noiseless settings, model well-specification, asymptotic approximations only holding locally). The authors should reflect on how these assumptions might be violated in practice and what the implications would be.
        \item The authors should reflect on the scope of the claims made, e.g., if the approach was only tested on a few datasets or with a few runs. In general, empirical results often depend on implicit assumptions, which should be articulated.
        \item The authors should reflect on the factors that influence the performance of the approach. For example, a facial recognition algorithm may perform poorly when image resolution is low or images are taken in low lighting. Or a speech-to-text system might not be used reliably to provide closed captions for online lectures because it fails to handle technical jargon.
        \item The authors should discuss the computational efficiency of the proposed algorithms and how they scale with dataset size.
        \item If applicable, the authors should discuss possible limitations of their approach to address problems of privacy and fairness.
        \item While the authors might fear that complete honesty about limitations might be used by reviewers as grounds for rejection, a worse outcome might be that reviewers discover limitations that aren't acknowledged in the paper. The authors should use their best judgment and recognize that individual actions in favor of transparency play an important role in developing norms that preserve the integrity of the community. Reviewers will be specifically instructed to not penalize honesty concerning limitations.
    \end{itemize}

\item {\bf Theory assumptions and proofs}
    \item[] Question: For each theoretical result, does the paper provide the full set of assumptions and a complete (and correct) proof?
    \item[] Answer: \answerYes{}
    \item[] Justification: Every theoretical result reported as established in the paper has a complete Lean proof in the supplementary artifact, and the relevant repository code compiles.  The theorem signatures expose the formal assumptions checked by Lean.  This includes Formalization~3 and the common-domain-to-trial-residual bridge, as well as the complete Proposition~4.4 pipeline: the printed proposition is represented formally, an explicit counterexample proves its negation, and the $Q$-norm repair is a separate checked theorem.
    \item[] Guidelines:
    \begin{itemize}
        \item The answer \answerNA{} means that the paper does not include theoretical results. 
        \item All the theorems, formulas, and proofs in the paper should be numbered and cross-referenced.
        \item All assumptions should be clearly stated or referenced in the statement of any theorems.
        \item The proofs can either appear in the main paper or the supplemental material, but if they appear in the supplemental material, the authors are encouraged to provide a short proof sketch to provide intuition. 
        \item Inversely, any informal proof provided in the core of the paper should be complemented by formal proofs provided in appendix or supplemental material.
        \item Theorems and Lemmas that the proof relies upon should be properly referenced. 
    \end{itemize}

    \item {\bf Experimental result reproducibility}
    \item[] Question: Does the paper fully disclose all the information needed to reproduce the main experimental results of the paper to the extent that it affects the main claims and/or conclusions of the paper (regardless of whether the code and data are provided or not)?
    \item[] Answer: \answerYes{}
    \item[] Justification: The source tables, generation scripts, historical witnesses, and formalization code needed for the reported case study are included in the supplementary material; public repository URLs are withheld from the anonymous PDF.
    \item[] Guidelines:
    \begin{itemize}
        \item The answer \answerNA{} means that the paper does not include experiments.
        \item If the paper includes experiments, a \answerNo{} answer to this question will not be perceived well by the reviewers: Making the paper reproducible is important, regardless of whether the code and data are provided or not.
        \item If the contribution is a dataset and\slash or model, the authors should describe the steps taken to make their results reproducible or verifiable. 
        \item Depending on the contribution, reproducibility can be accomplished in various ways. For example, if the contribution is a novel architecture, describing the architecture fully might suffice, or if the contribution is a specific model and empirical evaluation, it may be necessary to either make it possible for others to replicate the model with the same dataset, or provide access to the model. In general. releasing code and data is often one good way to accomplish this, but reproducibility can also be provided via detailed instructions for how to replicate the results, access to a hosted model (e.g., in the case of a large language model), releasing of a model checkpoint, or other means that are appropriate to the research performed.
        \item While NeurIPS does not require releasing code, the conference does require all submissions to provide some reasonable avenue for reproducibility, which may depend on the nature of the contribution. For example
        \begin{enumerate}
            \item If the contribution is primarily a new algorithm, the paper should make it clear how to reproduce that algorithm.
            \item If the contribution is primarily a new model architecture, the paper should describe the architecture clearly and fully.
            \item If the contribution is a new model (e.g., a large language model), then there should either be a way to access this model for reproducing the results or a way to reproduce the model (e.g., with an open-source dataset or instructions for how to construct the dataset).
            \item We recognize that reproducibility may be tricky in some cases, in which case authors are welcome to describe the particular way they provide for reproducibility. In the case of closed-source models, it may be that access to the model is limited in some way (e.g., to registered users), but it should be possible for other researchers to have some path to reproducing or verifying the results.
        \end{enumerate}
    \end{itemize}


\item {\bf Open access to data and code}
    \item[] Question: Does the paper provide open access to the data and code, with sufficient instructions to faithfully reproduce the main experimental results, as described in supplemental material?
    \item[] Answer: \answerYes{}
    \item[] Justification: The supplementary archive contains the formalization, evidence-generation scripts, structured review records, and build instructions used for the reported results and figures.
    \item[] Guidelines:
    \begin{itemize}
        \item The answer \answerNA{} means that paper does not include experiments requiring code.
        \item Please see the NeurIPS code and data submission guidelines (\url{https://neurips.cc/public/guides/CodeSubmissionPolicy}) for more details.
        \item While we encourage the release of code and data, we understand that this might not be possible, so \answerNo{} is an acceptable answer. Papers cannot be rejected simply for not including code, unless this is central to the contribution (e.g., for a new open-source benchmark).
        \item The instructions should contain the exact command and environment needed to run to reproduce the results. See the NeurIPS code and data submission guidelines (\url{https://neurips.cc/public/guides/CodeSubmissionPolicy}) for more details.
        \item The authors should provide instructions on data access and preparation, including how to access the raw data, preprocessed data, intermediate data, and generated data, etc.
        \item The authors should provide scripts to reproduce all experimental results for the new proposed method and baselines. If only a subset of experiments are reproducible, they should state which ones are omitted from the script and why.
        \item At submission time, to preserve anonymity, the authors should release anonymized versions (if applicable).
        \item Providing as much information as possible in supplemental material (appended to the paper) is recommended, but including URLs to data and code is permitted.
    \end{itemize}


\item {\bf Experimental setting/details}
    \item[] Question: Does the paper specify all the training and test details (e.g., data splits, hyperparameters, how they were chosen, type of optimizer) necessary to understand the results?
    \item[] Answer: \answerNA{}
    \item[] Justification: The paper does not report a controlled empirical experiment with tunable experimental settings.  Its quantitative material is descriptive project and literature-accounting evidence.
    \item[] Guidelines:
    \begin{itemize}
        \item The answer \answerNA{} means that the paper does not include experiments.
        \item The experimental setting should be presented in the core of the paper to a level of detail that is necessary to appreciate the results and make sense of them.
        \item The full details can be provided either with the code, in appendix, or as supplemental material.
    \end{itemize}

\item {\bf Experiment statistical significance}
    \item[] Question: Does the paper report error bars suitably and correctly defined or other appropriate information about the statistical significance of the experiments?
    \item[] Answer: \answerNA{}
    \item[] Justification: The paper does not report repeated stochastic experiments or statistical hypothesis tests.
    \item[] Guidelines:
    \begin{itemize}
        \item The answer \answerNA{} means that the paper does not include experiments.
        \item The authors should answer \answerYes{} if the results are accompanied by error bars, confidence intervals, or statistical significance tests, at least for the experiments that support the main claims of the paper.
        \item The factors of variability that the error bars are capturing should be clearly stated (for example, train/test split, initialization, random drawing of some parameter, or overall run with given experimental conditions).
        \item The method for calculating the error bars should be explained (closed form formula, call to a library function, bootstrap, etc.)
        \item The assumptions made should be given (e.g., Normally distributed errors).
        \item It should be clear whether the error bar is the standard deviation or the standard error of the mean.
        \item It is OK to report 1-sigma error bars, but one should state it. The authors should preferably report a 2-sigma error bar than state that they have a 96\% CI, if the hypothesis of Normality of errors is not verified.
        \item For asymmetric distributions, the authors should be careful not to show in tables or figures symmetric error bars that would yield results that are out of range (e.g., negative error rates).
        \item If error bars are reported in tables or plots, the authors should explain in the text how they were calculated and reference the corresponding figures or tables in the text.
    \end{itemize}

\item {\bf Experiments compute resources}
    \item[] Question: For each experiment, does the paper provide sufficient information on the computer resources (type of compute workers, memory, time of execution) needed to reproduce the experiments?
    \item[] Answer: \answerNA{}
    \item[] Justification: The paper does not present an experiment whose compute budget determines a reported performance result.  Appendix~\ref{app:resources} separately reports incomplete retained LLM-resource telemetry for the formalization effort.
    \item[] Guidelines:
    \begin{itemize}
        \item The answer \answerNA{} means that the paper does not include experiments.
        \item The paper should indicate the type of compute workers CPU or GPU, internal cluster, or cloud provider, including relevant memory and storage.
        \item The paper should provide the amount of compute required for each of the individual experimental runs as well as estimate the total compute. 
        \item The paper should disclose whether the full research project required more compute than the experiments reported in the paper (e.g., preliminary or failed experiments that didn't make it into the paper). 
    \end{itemize}
    
\item {\bf Code of ethics}
    \item[] Question: Does the research conducted in the paper conform, in every respect, with the NeurIPS Code of Ethics \url{https://neurips.cc/public/EthicsGuidelines}?
    \item[] Answer: \answerYes{}
    \item[] Justification: The work is a formal-methods case study using public mathematical sources and public first-person accounts; no recruited participants or private datasets are used.
    \item[] Guidelines:
    \begin{itemize}
        \item The answer \answerNA{} means that the authors have not reviewed the NeurIPS Code of Ethics.
        \item If the authors answer \answerNo, they should explain the special circumstances that require a deviation from the Code of Ethics.
        \item The authors should make sure to preserve anonymity (e.g., if there is a special consideration due to laws or regulations in their jurisdiction).
    \end{itemize}


\item {\bf Broader impacts}
    \item[] Question: Does the paper discuss both potential positive societal impacts and negative societal impacts of the work performed?
    \item[] Answer: \answerNA{}
    \item[] Justification: The work is foundational formal-methods research and does not study a deployed system or make a direct application claim requiring a separate societal-impact analysis.
    \item[] Guidelines:
    \begin{itemize}
        \item The answer \answerNA{} means that there is no societal impact of the work performed.
        \item If the authors answer \answerNA{} or \answerNo, they should explain why their work has no societal impact or why the paper does not address societal impact.
        \item Examples of negative societal impacts include potential malicious or unintended uses (e.g., disinformation, generating fake profiles, surveillance), fairness considerations (e.g., deployment of technologies that could make decisions that unfairly impact specific groups), privacy considerations, and security considerations.
        \item The conference expects that many papers will be foundational research and not tied to particular applications, let alone deployments. However, if there is a direct path to any negative applications, the authors should point it out. For example, it is legitimate to point out that an improvement in the quality of generative models could be used to generate Deepfakes for disinformation. On the other hand, it is not needed to point out that a generic algorithm for optimizing neural networks could enable people to train models that generate Deepfakes faster.
        \item The authors should consider possible harms that could arise when the technology is being used as intended and functioning correctly, harms that could arise when the technology is being used as intended but gives incorrect results, and harms following from (intentional or unintentional) misuse of the technology.
        \item If there are negative societal impacts, the authors could also discuss possible mitigation strategies (e.g., gated release of models, providing defenses in addition to attacks, mechanisms for monitoring misuse, mechanisms to monitor how a system learns from feedback over time, improving the efficiency and accessibility of ML).
    \end{itemize}
    
\item {\bf Safeguards}
    \item[] Question: Does the paper describe safeguards that have been put in place for responsible release of data or models that have a high risk for misuse (e.g., pre-trained language models, image generators, or scraped datasets)?
    \item[] Answer: \answerNA{}
    \item[] Justification: The paper does not release a high-risk model or dataset requiring misuse safeguards.
    \item[] Guidelines:
    \begin{itemize}
        \item The answer \answerNA{} means that the paper poses no such risks.
        \item Released models that have a high risk for misuse or dual-use should be released with necessary safeguards to allow for controlled use of the model, for example by requiring that users adhere to usage guidelines or restrictions to access the model or implementing safety filters. 
        \item Datasets that have been scraped from the Internet could pose safety risks. The authors should describe how they avoided releasing unsafe images.
        \item We recognize that providing effective safeguards is challenging, and many papers do not require this, but we encourage authors to take this into account and make a best faith effort.
    \end{itemize}

\item {\bf Licenses for existing assets}
    \item[] Question: Are the creators or original owners of assets (e.g., code, data, models), used in the paper, properly credited and are the license and terms of use explicitly mentioned and properly respected?
    \item[] Answer: \answerYes{}
    \item[] Justification: The manuscript cites the mathematical, software, and public-source assets used in the argument.  The supplementary code retains the upstream package and repository provenance used by the formalization.
    \item[] Guidelines:
    \begin{itemize}
        \item The answer \answerNA{} means that the paper does not use existing assets.
        \item The authors should cite the original paper that produced the code package or dataset.
        \item The authors should state which version of the asset is used and, if possible, include a URL.
        \item The name of the license (e.g., CC-BY 4.0) should be included for each asset.
        \item For scraped data from a particular source (e.g., website), the copyright and terms of service of that source should be provided.
        \item If assets are released, the license, copyright information, and terms of use in the package should be provided. For popular datasets, \url{paperswithcode.com/datasets} has curated licenses for some datasets. Their licensing guide can help determine the license of a dataset.
        \item For existing datasets that are re-packaged, both the original license and the license of the derived asset (if it has changed) should be provided.
        \item If this information is not available online, the authors are encouraged to reach out to the asset's creators.
    \end{itemize}

\item {\bf New assets}
    \item[] Question: Are new assets introduced in the paper well documented and is the documentation provided alongside the assets?
    \item[] Answer: \answerYes{}
    \item[] Justification: The supplementary material documents the project code, source records, generated evidence, and build commands in its README files.
    \item[] Guidelines:
    \begin{itemize}
        \item The answer \answerNA{} means that the paper does not release new assets.
        \item Researchers should communicate the details of the dataset\slash code\slash model as part of their submissions via structured templates. This includes details about training, license, limitations, etc. 
        \item The paper should discuss whether and how consent was obtained from people whose asset is used.
        \item At submission time, remember to anonymize your assets (if applicable). You can either create an anonymized URL or include an anonymized zip file.
    \end{itemize}

\item {\bf Crowdsourcing and research with human subjects}
    \item[] Question: For crowdsourcing experiments and research with human subjects, does the paper include the full text of instructions given to participants and screenshots, if applicable, as well as details about compensation (if any)? 
    \item[] Answer: \answerNA{}
    \item[] Justification: The paper uses published public first-person accounts as documentary sources and does not recruit, instruct, or compensate participants.
    \item[] Guidelines:
    \begin{itemize}
        \item The answer \answerNA{} means that the paper does not involve crowdsourcing nor research with human subjects.
        \item Including this information in the supplemental material is fine, but if the main contribution of the paper involves human subjects, then as much detail as possible should be included in the main paper. 
        \item According to the NeurIPS Code of Ethics, workers involved in data collection, curation, or other labor should be paid at least the minimum wage in the country of the data collector. 
    \end{itemize}

\item {\bf Institutional review board (IRB) approvals or equivalent for research with human subjects}
    \item[] Question: Does the paper describe potential risks incurred by study participants, whether such risks were disclosed to the subjects, and whether Institutional Review Board (IRB) approvals (or an equivalent approval/review based on the requirements of your country or institution) were obtained?
    \item[] Answer: \answerNA{}
    \item[] Justification: No crowdsourcing or recruited human-subject experiment is conducted; the practitioner accounts are previously published public documents.
    \item[] Guidelines:
    \begin{itemize}
        \item The answer \answerNA{} means that the paper does not involve crowdsourcing nor research with human subjects.
        \item Depending on the country in which research is conducted, IRB approval (or equivalent) may be required for any human subjects research. If you obtained IRB approval, you should clearly state this in the paper. 
        \item We recognize that the procedures for this may vary significantly between institutions and locations, and we expect authors to adhere to the NeurIPS Code of Ethics and the guidelines for their institution. 
        \item For initial submissions, do not include any information that would break anonymity (if applicable), such as the institution conducting the review.
    \end{itemize}

\item {\bf Declaration of LLM usage}
    \item[] Question: Does the paper describe the usage of LLMs if it is an important, original, or non-standard component of the core methods in this research? Note that if the LLM is used only for writing, editing, or formatting purposes and does \emph{not} impact the core methodology, scientific rigor, or originality of the research, declaration is not required.
    %this research? 
    \item[] Answer: \answerYes{}
    \item[] Justification: The main text identifies the LLM systems and their roles, and Appendix~\ref{app:resources} reports the model families and the recovered resource telemetry available for the project.
    \item[] Guidelines:
    \begin{itemize}
        \item The answer \answerNA{} means that the core method development in this research does not involve LLMs as any important, original, or non-standard components.
        \item Please refer to our LLM policy in the NeurIPS handbook for what should or should not be described.
    \end{itemize}

\end{enumerate}

\end{document}
